\subsection{Кодовые фрагменты, иллюстрирующие реализацию}

Для подтверждения того, что изложенные проектные решения реализованы в кодовой
базе, в работе приведены листинги ключевых фрагментов программной системы.
Они выбраны не по признаку краткости, а по связи с критическими требованиями:
транзакционная согласованность, идемпотентность обработки и управляемость
жизненного цикла заданий индексирования.

На рисунке~\ref{src:src1} приведен фрагмент метода сохранения блока. Он
показывает начало транзакционной записи, использование advisory lock для
синхронизации состояния цепи, проверку повторной обработки уже сохраненной
высоты и ожидание предыдущего блока перед записью следующей высоты. Данный
фрагмент непосредственно связан с требованиями корректной обработки
подтвержденной цепи, идемпотентности и устойчивости к повторному запуску job.

\begin{figure}
\lstinputlisting[firstline=77,lastline=96]{../src/modules/indexer/mod.rs}
\caption{Фрагмент транзакционной и идемпотентной записи блока}
\label{src:src1}
\end{figure}

На рисунке~\ref{src:src2} приведена функция, задающая допустимые переходы
состояний job. Именно этот фрагмент делает жизненный цикл задания
формализованным контрактом: недопустимые переходы возвращают ошибку, а REST API,
CLI и административная панель используют единую бизнес-логику управления
индексированием.

\begin{figure}
\lstinputlisting[firstline=575,lastline=585]{../src/modules/jobs/mod.rs}
\caption{Фрагмент проверки допустимых переходов job}
\label{src:src2}
\end{figure}

Командный интерфейс реализован в отдельном модуле \texttt{cli/indexer\_cli.py}
и используется для вспомогательного взаимодействия с заданиями индексирования,
узлами и прикладными данными через тот же REST API, что и административная
панель.
